% ═══════════════════════════════════════════════════════════════════
% Chapter 7 — Tensor IR: IR001
% ═══════════════════════════════════════════════════════════════════

\section{Tensor IR: \IR{1}}
\label{sec:ir001}

\IR{1} is Crucible's tensor-level intermediate representation: the form in which a captured machine-learning computation is held between Vessel dispatch capture (Chapter~4) and Forge's lowering pass (Chapter~8).
It is the highest tier of Crucible's three-tier IR stack, sitting above Forge's portable kernel IR \IR{2} (Chapter~9) and above the per-vendor machine IRs \IRstar{3} (Chapter~11).
\IR{1} is constructed on Crucible's background thread immediately after a TraceRing drain, and is the input to every Forge phase prior to phase~E (\textsc{lower\_to\_kernels}).

\IR{1} is vendor-neutral by construction.
It encodes operator identity, tensor shape, dataflow structure, and symbolic dimension algebra; it does not encode tile shapes, register allocations, kernel templates, address spaces, instruction-scheduling hints, or any other hardware-specific concept.
The boundary is structural: the \IR{1} type system exposes no field that could carry such information.
Adding hardware concerns to \IR{1} would require visible changes to its types, not merely to its content.
The deliberate boundary keeps the Vessel adapter layer (one per frontend) thin --- a Vessel knows how to translate its host framework's dispatch into \IR{1}, nothing more --- and concentrates all hardware adaptation into the Forge--Mimic compilation pipeline.

\subsection{Components}
\label{sec:ir001-components}

\IR{1} comprises five interacting components.

\textbf{Graph IR.}
The mutable computation graph (header \code{Graph.h}).
Each \struct{GraphNode} occupies exactly 64 bytes (one cache line) and carries node identity, kind, dimensionality, dtype, reduction operator, input-count, symbolic size and stride pointers, body or extern-info pointer, input-array pointer, use count, schedule order, group hash, and fused-group identifier.
Compute bodies are flat instruction arrays of 8-byte SSA micro-ops (\struct{Inst}) drawn from a fixed micro-op set (approximately fifty opcodes).
The Graph IR provides standard mutation primitives --- dead-code elimination via fixpoint, common-subexpression elimination via structural hashing, topological sort via Kahn's algorithm, and replace-all-uses-with --- all operating in $O(V{+}E)$ on counting-sort-built CSR adjacency.

\textbf{Symbolic expression algebra.}
\struct{Expr} (header \code{Expr.h}, pool \code{ExprPool.h}) is an interned, immutable expression algebra.
Each \struct{Expr} node is 32 bytes and carries the operator, operands, dtype, and a vector of assumption flags (\code{is\_integer}, \code{is\_real}, \code{is\_finite}, \code{is\_positive}, \code{is\_negative}, \code{is\_nonnegative}, \code{is\_nonpositive}, \code{is\_zero}, \code{is\_even}, \code{is\_odd}, \code{is\_number}, \code{is\_symbol}, \code{is\_boolean}).
The algebra covers arithmetic (add, subtract, multiply, divide), relational and logical operations, modular and floor-and-ceiling division, modulo, rounding, shifts, bitwise operations, conditionals (\code{where}, \code{min}, \code{max}), trigonometric and transcendental functions, negation, and bit reinterpretation (\code{bit\_cast}).
Expressions are interned via a Swiss-table over the wyhash function: structurally equal expressions share a pointer.
Two \struct{Expr*} compare equal if and only if they denote the same value --- structural equality reduces to pointer equality, the substrate primitive that drives all downstream symbolic analysis.

\textbf{Operator taxonomy.}
The \struct{CKernelId} catalog (header \code{CKernel.h}) is the device-agnostic operator namespace.
Every captured operation resolves either to a taxonomy entry --- specific identities for linear algebra, convolution, attention, normalization, activations, elementwise binary and unary operations, reductions, pooling, data movement, embedding, copy, vision, fused variants, linear-algebra decompositions, state-space models, production inference operations (paged attention, MoE routing, ragged attention), three-dimensional rendering, structured matrix and graph operations, collective communication, input--output, random-number generation, and synchronization --- or to an \code{OPAQUE} placeholder for genuinely unrecognized operations.
Taxonomy ordinals are frozen across releases for the substrate operators (Section~1 of the catalog); extensions land in a growth zone (Section~2) with strictly increasing ordinals.
Two-phase registration ensures coverage: each Vessel registers a \struct{SchemaHash}\,$\to$\,\struct{CKernelId} mapping at startup, and \code{classify\_kernel} resolves each captured operation during \code{build\_trace}.

\textbf{Symbol table.}
\struct{SymbolTable} (header \code{SymbolTable.h}) holds per-symbol metadata: kind, hint, declared range, and flags.
Symbols name the dynamic quantities that the symbolic algebra references --- typically dimensions of dynamic tensors (sequence length, batch size, vocabulary size), strides of non-contiguous tensors, and other scalar parameters bound at runtime.
Range information is propagated through the algebra and consumed downstream by Forge's bucketed shape specialization (Phase~F.4).

\textbf{Merkle structure.}
The content-addressed graph of compilable units lives in \code{MerkleDag.h}.
Three node kinds are defined:

\begin{itemize}
  \item \struct{RegionNode} --- a compilable operation sequence, the atomic unit of the kernel cache.
        A region carries an arena-allocated array of \struct{TraceEntry} records, a \struct{ContentHash} computed by wymix-folding the schema hashes, per-tensor dimension and stride products, packed dtype and device metadata, and scalar arguments;
        an atomic \struct{CompiledKernel*} pointer (background-written with release, foreground-read with acquire);
        a \struct{MemoryPlan*}; and provenance fields.
  \item \struct{BranchNode} --- a guard-based dispatch with a sorted array of arms for $O(\log n)$ lookup.
        Each arm is a \code{(predicate, target)} pair.
        The guard predicate is one of five kinds: shape-dimension mismatch, scalar-value mismatch, dtype mismatch, device mismatch, or operation-sequence mismatch.
        BranchNodes are the substrate for every form of runtime adaptation: dynamic-shape dispatch, architecture mutation, attention-head replacement, hyperparameter changes, and continuous learning with A/B verification.
  \item \struct{LoopNode} --- cyclic computation within the acyclic DAG framework.
        A LoopNode wraps an acyclic body sub-DAG with an array of feedback edges and a termination predicate (fixed iteration count or convergence within an epsilon).
        The Merkle hash of a LoopNode is salted to distinguish it from a RegionNode of the same body content.
        LoopNodes encode bounded recurrence (recurrent neural networks executing as one body $\times$ $N$ iterations without a Python step), convergence-based execution (DEQ fixed-point iteration, diffusion denoising), and cross-iteration pipelining.
\end{itemize}

Every Merkle node carries a \struct{MerkleHash} that recursively combines its content hash with the merkle hashes of its children, providing $O(1)$ subtree equality and $O(\log N)$ structural diff.

The five components compose: Graph IR provides the local optimization surface; \struct{Expr} provides the symbolic dimension and stride algebra; the operator taxonomy provides the unifying operator identity; the symbol table tracks dimensional unknowns; and the Merkle structure provides versioning, equality, and the kernel-cache key.

\subsection{Properties}
\label{sec:ir001-properties}

The properties below are commitments of the \IR{1} specification.
Each is enforced either structurally (by the type system or by component invariants) or by the construction algorithm (Section~\ref{sec:ir001-construction}).

\begin{property}[Vendor neutrality]
\label{prop:ir001-vendor-neutral}
\IR{1} contains no vendor-specific opcode, no vendor-specific layout, no vendor address-space tag, and no vendor-specific datatype beyond the \struct{ScalarType} enumeration.
Any well-formed \IR{1} graph is the input of every supported Forge configuration; Forge's choice of recipe and tile specialization (Phase~E) is independent of which Vessel produced the graph.
\end{property}

\begin{property}[Deterministic construction]
\label{prop:ir001-determinism}
Given the same TraceRing entries and the same MetaLog entries in the same order, \IR{1} construction produces a byte-identical graph.
Every node's \struct{ContentHash} is determined by the operator identity, input shapes, strides, dtypes, devices, and scalar argument values --- quantities invariant under replay.
This is the substrate of Crucible's replay-determinism guarantee (DetSafe, Chapter~22).
\end{property}

\begin{property}[Content addressing]
\label{prop:ir001-content-addressing}
Every \struct{RegionNode} carries a 64-bit \struct{ContentHash}.
Identical computations produce identical hashes --- across iterations, across runs, across models, and across organizations --- with collision probability bounded by the wymix mixer's avalanche distribution.
Every Merkle node carries a 64-bit \struct{MerkleHash} extending its content hash with the merkle hashes of its children, providing constant-time subtree equality.
\end{property}

\begin{property}[Affine-faithful symbolic algebra]
\label{prop:ir001-expr-canonical}
\struct{Expr} is a strongly normalizing affine-expression engine.
Any expression involving constants, symbols, and the supported operators reduces to a unique canonical form during interning.
Two expressions denoting the same value are interned to the same \struct{Expr*}; structural equality is therefore decidable in $O(1)$ via pointer comparison.
This is the substrate that downstream phases use for shape arithmetic, stride computation, loop-bound analysis, and address calculation.
\end{property}

\begin{property}[Closure under capture]
\label{prop:ir001-closure}
Every operator dispatched by the Vessel layer resolves either to a \struct{CKernelId} taxonomy entry or to the \code{OPAQUE} placeholder.
The taxonomy is closed over the machine-learning operator surface; \code{OPAQUE} provides the escape valve for genuinely unrecognized operators without requiring users to extend the IR.
\end{property}

\begin{property}[Compositionality with the safety substrate]
\label{prop:ir001-substrate-compose}
Arena allocation, non-static data-member initialization on every field, strong-typed identifiers (\struct{OpIndex}, \struct{SlotId}, \struct{NodeId}, \struct{SymbolId}, \struct{MetaIndex}), strong-typed hashes (\struct{SchemaHash}, \struct{ShapeHash}, \struct{ContentHash}, \struct{MerkleHash}, \struct{ScopeHash}, \struct{CallsiteHash}), and the eight safety axioms (Chapter~22) apply to every \IR{1} component without exception.
There is no fast path that bypasses the discipline.
\end{property}

\subsection{Construction algorithm}
\label{sec:ir001-construction}

\IR{1} is constructed on Crucible's background thread by \code{BackgroundThread::run}, which drains the TraceRing in batches and constructs the graph incrementally.
The procedure is single-pass: each ring entry is processed exactly once.
There is no backwards traversal, no second normalization pass, no fixpoint iteration.

\begin{enumerate}
  \item Drain up to \code{BATCH\_SIZE} entries from the TraceRing into a per-iteration arena buffer.
  \item For each entry, allocate a \struct{TraceEntry} record from the arena, populate operator identity, schema hash, shape hash, and copy \struct{TensorMeta} references from the MetaLog by index.
  \item Track tensor identity via a generation-counter \struct{PtrMap} (open-addressing, no per-iteration \code{memset}): a match between an operation's input \code{data\_ptr} and a previously recorded output \code{data\_ptr} produces a \code{DATA\_FLOW} edge; aliased pointers (views, in-place mutations) produce \code{ALIAS} edges.
  \item Construct bidirectional CSR adjacency simultaneously via prefix-sum and scatter.
        Both source-sorted and destination-sorted adjacencies are built in $O(V{+}E)$ in a single sweep.
  \item Detect iteration boundaries via the \struct{IterationDetector}: a $K{=}5$ schema-hash signature with two-match confirmation.
        The first signature match locks a candidate iteration boundary; the second match confirms the boundary.
        False-positive probability is bounded by $5 \times 2^{-320}$ under the hash-independence assumption.
  \item At each confirmed boundary, fold the accumulated \struct{TraceEntry} array into a \struct{RegionNode} by computing its \struct{ContentHash} (wymix fold) and recursively its \struct{MerkleHash}.
  \item Assign the new \struct{RegionNode} to the active graph chain via release-store on an atomic pointer, and signal Forge's compilation pipeline.
\end{enumerate}

The cost of construction is dominated by the CSR build (linear in vertex and edge count) and the content-hash fold (linear in the entry size).
On a representative thousand-operation graph, \IR{1} construction completes in approximately 50 to 100 microseconds.
All allocation is from the per-iteration arena; the arena is reset at iteration boundary, providing bulk free at $O(1)$.

\subsection{What \IR{1} does not contain}
\label{sec:ir001-not}

The boundary of \IR{1} is as load-bearing as its content.
\IR{1} does not contain any of the following.

\begin{itemize}
  \item Tile shapes, MMA-instruction shapes, vectorization widths, register counts, occupancy targets, or any other hardware-level resource decision.
        These appear at \IR{2} (Chapter~9) and are committed by Forge's Phase~F.
  \item \struct{NumericalRecipe} choices (accumulator dtype, reduction algorithm, rounding mode, scale policy, softmax recurrence variant, determinism tier).
        These are pinned at \IR{2} by Forge's Phase~E (\textsc{lower\_to\_kernels}).
  \item Address-space tags (global memory, shared memory, tensor memory, distributed shared memory), register-bank assignments, scoreboard slots, instruction-scheduling hints, or pipeline depths.
        These appear at \IRstar{3} (Chapter~11) and are committed by the per-vendor Mimic backend.
  \item Vendor library calls (cuBLAS, cuDNN, cuSPARSE, NCCL, rocBLAS, MIOpen, RCCL, libtpu, libnrt, libnccom, hcoll, libsharp), vendor compiler intermediates (PTX, HLO, NEFF, SPIR-V), or vendor binary format wrappers (cubin, HSACO, TPU executable, NEFF tarball).
        These never appear in any Crucible IR; per-vendor Mimic backends emit the binary format directly from \IRstar{3}.
\end{itemize}

The discipline is structural.
The \IR{1} type system exposes no field that could carry the prohibited information; adding such information would require visible type changes, not merely content changes, and would be caught at the boundary by inspection or by the negative-compile fixtures of the safety substrate (Chapter~22).

\subsection{Relationship to downstream IRs}
\label{sec:ir001-downstream}

The compilation pipeline consumes \IR{1} in stages.
Forge's Phases~A through D (Chapter~8) operate on \IR{1} in place: canonicalization, type and shape analysis, exhaustive rewriting, and global fusion.
Phase~E lowers each \struct{FusedRegion} to one or more \struct{KernelNode}s in \IR{2}; this is the boundary at which kernel templates are matched, semantic layouts are committed, and \struct{NumericalRecipe}s are pinned.
Phases~F and G operate on \IR{2}; Phase~H delegates each \struct{KernelNode} to the appropriate Mimic per-vendor backend, which lowers \IR{2} to its specific \IRstar{3} variant, performs MAP-Elites kernel search guided by a calibrated three-tier simulator, and emits the native binary format.

\IR{1}, as the highest tier, sees only what is invariant across this pipeline: the operation identity, the tensor metadata, the symbolic shape and stride algebra, and the Merkle structure that anchors the kernel cache.
Every choice that depends on the target hardware is deferred to \IR{2} and \IRstar{3}.
The deferral is what makes the kernel cache federation-safe at the L1 layer (Chapter~21): two installations running the same model produce byte-identical \IR{1} and byte-identical \IR{2} (when their fleet recipe-intersections agree), regardless of the specific hardware they target.
